\title{xLSTM: Extended Long Short-Term Memory}

\begin{abstract}
During the 1990s, the introduction of the constant error carousel and gating mechanisms established the foundational principles of Long Short-Term Memory (LSTM). Since their inception, LSTMs have proven to be robust and instrumental in a wide array of deep learning achievements, notably serving as the architecture behind the first Large Language Models (LLMs). The emergence of Transformer architectures, driven by parallelizable self-attention, has, however, shifted the paradigm and surpassed LSTMs in large-scale applications. This leads us to ask: what can be achieved in language modeling by extending LSTMs to the scale of billions of parameters, employing contemporary techniques from modern LLMs, while addressing LSTMs' traditional weaknesses? To explore this, we propose the use of exponential gating combined with normalization and stabilization strategies. Furthermore, we reconfigure the LSTM memory, introducing: (i) sLSTM, which employs scalar memory and update mechanisms with an updated mixing approach, and (ii) mLSTM, which is designed for complete parallelization through matrix memory and a covariance-based update scheme. By embedding these LSTM enhancements within residual block frameworks, we form xLSTM blocks, which are further composed residually into xLSTM-based architectures. The incorporation of exponential gating and revised memory configurations significantly enhances the performance of xLSTM, enabling competitive results relative to state-of-the-art Transformers and State Space Models in terms of both efficacy and scalability.\\
Code available at: \url{https://github.com/NX-AI/xlstm}
\end{abstract}

\section{Introduction}
\label{sec:intro}

The concepts underlying Long Short-Term Memory (LSTM)~\citep{Hochreiter:91,Hochreiter:97nips,Hochreiter:97}, specifically the constant error carousel and the mechanism of gating, were proposed as a solution to the vanishing gradient issue encountered in recurrent neural networks~\citep{Hochreiter:91,Hochreiter:00book}:

\begin{align}
\label{eq:lstm_idea}
  c_t \ = \ f_t \ c_{t-1} \ + \ 
          i_t \  z_t \ , \quad  
          h_t \ = \ o_t \ \psi({c_t}) \ . 
\end{align}

The constant error carousel refers to the process where the previous cell state $c_{t-1}$ (shown in green) receives an additive modification from the cell inputs $z_t$, with this interaction being regulated by sigmoid gates (indicated in blue). The input gate $i_t$ and forget gate $f_t$ serve to govern this update to the cell state, whereas the output gate $o_t$ determines the release of information from the memory cell, specifically producing the hidden state $h_t$. The cell state undergoes normalization or compression via $\psi$ before the output gating mechanism generates the hidden state.

LSTMs have achieved notable success across multiple fields \citep{Hochreiter:01,Hochreiter:07,Schmidhuber:15}, maintaining dominance in text generation up until the introduction of Transformers in 2017~\citep{Vaswani:17}. Their capabilities have been validated in a wide range of sequential tasks, including but not limited to, text synthesis~\citep{Graves:13,Karpathy:15blog}, handwriting generation~\citep{Graves:13}, sequence-to-sequence translation~\citep{Sutskever:14nips}, program evaluation~\citep{Zaremba:14arxiv}, automated image captioning~\citep{Karpathy:15,Hossain:19}, source code generation~\citep{Karpathy:15blog}, rainfall-runoff simulations~\citep{Kratzert:18,Kratzert:19}, and hydrological modeling for flood prediction~\citep{Nearing:24}. In the context of reinforcement learning, LSTMs remain the leading sequence modeling architecture, as exemplified by their integration in the AlphaStar agent for StarCraft II~\citep{Vinyals:17short}, OpenAI Five for Dota 2~\citep{OpenAI:19}, and magnetic controller models for nuclear fusion reactors~\citep{Degrave:22short}. Their proficiency in abstract representation—manifested through their ability to encode and store semantic features within memory cells~\citep{Karpathy:15blog}—has been illustrated by the emergence of specialized neurons encoding number and syntax~\citep{Lakretz:19}, linguistic features~\citep{Bau:19}, and sentiment~\citep{Radford:17arxiv}. LSTMs continue to find application in crucial, contemporary systems~\citep{Degrave:22short,Nearing:24}, underscoring their enduring relevance.

Although LSTMs have achieved remarkable results, they are hampered by three principal shortcomings: (i) Inability to update previously made storage choices. We illustrate this drawback with the {\it Nearest Neighbor Search} task (refer to Appendix~\ref{sec:appNearestNeighborSearch}): Given a reference vector, the system must process a sequence element-by-element to find and retrieve the value associated with the most similar vector. As depicted in the left panel of Figure~\ref{fig:lstmProblems}, LSTMs exhibit higher mean squared error on this task, reflecting difficulty in overwriting stored information when a closer match is later encountered. In contrast, our proposed xLSTM addresses this issue using exponential gating. (ii) Constrained capacity for information retention, as traditional LSTMs must condense memory into one-dimensional cell states. We demonstrate this limitation using the {\it Rare Token Prediction} task. The right panel of Figure~\ref{fig:lstmProblems} displays token prediction perplexity on Wikitext-103~\citep{Merity:17}, segmented by token frequency. LSTMs underperform on infrequent tokens, underscoring their restricted memory capacity, while the xLSTM overcomes this constraint through matrix-based memory. (iii) Poor parallelization resulting from sequential state dependencies, due to the recurrence that couples hidden states across time steps, thereby necessitating strictly sequential computation.

The shortcomings of LSTM have facilitated the rise of Transformers~\citep{Vaswani:17} within the field of language modeling. An intriguing question arises: if we address these limitations and expand LSTMs to match the scale of today's Large Language Models, what levels of language modeling performance are attainable?

\section{Extended Long Short-Term Memory}
\label{sec:xLSTM}

In response to the inherent constraints of LSTM, the Extended Long Short-Term Memory (xLSTM) framework incorporates two pivotal changes to the foundational LSTM concept presented in Equation~\eqref{eq:lstm_idea}. These innovations—namely exponential gating and advanced memory structures—lead to two novel variants within the LSTM family: (i) the sLSTM (discussed in Section~\ref{sec:sLSTM}), characterized by its scalar memory, scalar update mechanism, and incorporation of memory mixing; and (ii) the mLSTM (outlined in Section~\ref{sec:mLSTM}), which utilizes a matrix-based memory representation alongside an update mechanism based on the covariance (outer product), designed for complete parallelizability. Both of these architectures leverage exponential gating to advance beyond standard LSTM performance. To achieve efficient parallel computation, the mLSTM forgoes the use of memory mixing and the corresponding hidden-hidden recurrent links. Additionally, the sLSTM and mLSTM models can be configured with multiple memory cells. In such configurations, sLSTM supports inter-cell memory mixing, whereas sLSTM can further extend to multiple attention heads, where memory mixing is confined to within each head among its respective cells and does not occur between heads. This arrangement—combining exponential gating with multi-headed sLSTM—yields a novel methodology for memory mixing. For the mLSTM architecture, the distinction between multiple heads and multiple cells is effectively eliminated, as both are functionally identical.

By embedding the novel LSTM variants within residual block modules, we obtain xLSTM blocks (refer to Section~\ref{sec:xLSTMblock}). Arranging these xLSTM blocks in a residual stacking manner gives rise to xLSTM architectures, which are discussed in Section~\ref{sec:xLSTMarchitecure}. The composition of the xLSTM architecture and its integral elements is illustrated in Figure~\ref{fig:xlstm_sketch}.

\subsection{Review of the Long Short-Term Memory}
\label{sec:orgLSTM}

The initial concept of LSTM~\citep{Hochreiter:91,Hochreiter:97nips,Hochreiter:97} proposed a scalar memory cell serving as the main unit for computation and storage, which circumvents the vanishing gradient issue~\citep{Hochreiter:91,Hochreiter:00book} via the constant error carousel mechanism responsible for updating the cell state. This memory cell incorporates three gates: the input gate, output gate, and forget gate, with the latter being introduced in the work of \citet{Gers:00}. The following equations define how the LSTM memory cell is updated at each time step $t$:

\begin{align}
c_t \ &= \  f_t \ c_{t-1} \ + \ i_t \ z_t &  & &\text{cell state} \\
h_t  \ &= \ o_t \ \tilde{h}_t \ , & \tilde{h}_t \ &= \ \psi \left( c_t\right)
  &\text{hidden state} \\
z_t \ &= \ \varphi \left( \tilde{z}_t \right) \ , 
  &\tilde{z}_t \ &=  \ \mathbf{w}^\top_{z} \ \mathbf{x}_t \ + \
  r_{z}  h_{t-1} \ + \  b_{z} \ \
  &\text{cell input} \\
i_t \ &= \ \sigma \left( \tilde{i}_t  \right) \ , 
  &\tilde{i}_t \ &= \ \mathbf{w}^\top_{i} \ \mathbf{x}_t \ + \
  r_{i}  \ h_{t-1} \ + \  b_{i} \ \
  &\text{input gate} \\
f_t \ &= \ \sigma \left(  \tilde{f}_t \right) \ , 
  &\tilde{f}_t \ &= \ \mathbf{w}^\top_{f} \ \mathbf{x}_t  \ + \
  r_{f}  \ h_{t-1} \ + \  b_{f} \ \
  &\text{forget gate} \\
o_t \ &= \ \sigma \left( \tilde{o}_t \right) \ , 
  &\tilde{o}_t  \ &= \ \mathbf{w}^\top_{o} \ \mathbf{x}_t \ + \
  r_{o}  \ h_{t-1} \ + \  b_{o} \ \
  &\text{output gate} 
\end{align}

The input weight vectors $\mathbf{w}_{z}$, $\mathbf{w}_{i}$, $\mathbf{w}_{f}$, and $\mathbf{w}_{o}$ denote the connections from the input $\mathbf{x}_t$ to the cell input, input gate, forget gate, and output gate, respectively. The recurrent weight terms $r_{z}$, $r_{i}$, $r_{f}$, and $r_{o}$ represent the links from the previous hidden state $h_{t-1}$ to each of these gates and the cell input. Associated bias parameters are $b_{z}$, $b_{i}$, $b_{f}$, and $b_{o}$. The activation functions $\varphi$ and $\psi$ (commonly $\tanh$) are used for the cell input and the hidden state; specifically, $\psi$ is applied to constrain the range of the cell state, preventing it from diverging. Sigmoid functions, defined as $\sigma \left( x \right)= 1/(1+\exp(-x))$, are applied to all gates. In subsequent extensions, collections of scalar memory cells $\mathbf{c}_t \in \mathbb{R}^{d}$ were grouped into vectors, enabling the implementation of recurrent weight matrices $\mathbf{R} \in \mathbb{R}^{d \times d}$ that facilitate interaction among memory cell outputs~\citep{Greff:15}. Further details are provided in Appendix~\ref{sec:apporgLSTM}. Ablation analyses have demonstrated that each element within the memory cell architecture serves an essential function~\citep{Greff:15}.

\subsection{sLSTM}
\label{sec:sLSTM}

To enable LSTMs to modify storage decisions dynamically, we propose the use of exponential gates (highlighted in red) alongside mechanisms for normalization and stabilization. Specifically, we allow the input and forget gates to utilize exponential activation functions. For normalization, we add a dedicated normalizer state, which aggregates the product of the input gate and every subsequent forget gate over time.

The forward pass for the scalar sLSTM is as follows:

\begin{align}
\label{eq:slstmforward}
c_t \ &= \  f_t \ c_{t-1} \ + \ i_t \ z_t & & 
 &\text{cell state} \\
n_t \ &= \  f_t \ n_{t-1} \ + \ i_t  & & 
 &\text{normalizer state} \\
h_t \ &= \ o_t \ \tilde{h}_t \ ,
  &\tilde{h}_t \ &= \ c_t / n_t
&\text{hidden state} \\
z_t \ &= \ \varphi \left( \tilde{z}_t \right) \ , 
  &\tilde{z}_t \ &=  \ \mathbf{w}^\top_{z} \ \mathbf{x}_t \ + \
  r_{z}  h_{t-1} \ + \  b_{z}
  &\text{cell input} \\
i_t \ &= \ \!\exp \left( \tilde{i}_t  \right) \ , 
  &\tilde{i}_t \ &= \ \mathbf{w}^\top_{i} \ \mathbf{x}_t \ + \
  r_{i}  \ h_{t-1} \ + \  b_{i}
  &\text{input gate} \\
f_t \ &= \ \sigma \left(  \tilde{f}_t \right) \ \text{OR} \
\!\exp \left(  \tilde{f}_t \right) \ ,
  &\tilde{f}_t \ &= \ \mathbf{w}^\top_{f} \ \mathbf{x}_t  \ + \
  r_{f}  \ h_{t-1} \ + \  b_{f}
  &\text{forget gate} \\
o_t \ &= \ \sigma \left( \tilde{o}_t \right) \ , 
  &\tilde{o}_t  \ &= \ \mathbf{w}^\top_{o} \ \mathbf{x}_t \ + \
  r_{o}  \ h_{t-1} \ + \  b_{o}
  &\text{output gate} 
\end{align}

The gating mechanisms from the original LSTM, including input- and hidden-dependent gates along with the bias term, are incorporated into the redesigned architectures. Since exponential activation functions are prone to producing excessively large outputs that may result in numerical overflow, we incorporate an extra state \mbox{$m_t$~\citep{Milakov:18arxiv}} to help stabilize the gating process.

\begin{align}
\label{eq:slstmstabil}
m_t \ &= \ \max \left( \log ( f_t ) + m_{t-1} , \log ( i_t ) \right) &\text{stabilizer state} \\
i'_t \ &= \ \!\exp \left( \log \left ( i_t \right) - m_t \right) \ = \ \!\exp \left( \tilde{i}_t - m_t \right) \ \, 
  &\text{stabil. input gate} \\
f'_t \ &= \ \!\exp \left( \log \left( f_t \right) + m_{t-1} - m_t \right) \ \, 
  &\text{stabil. forget gate}
\end{align}

As demonstrated in Appendix~\ref{sec:appsLSTM}, substituting $f_t$ with $f'_t$ and $i_t$ with $i'_t$ during the forward computation leaves both the overall network output and the loss derivatives with respect to the parameters unaffected.

\paragraph{New Memory Mixing.} Like its predecessor, sLSTM supports several memory cells (refer to Appendix~\ref{sec:appsLSTM} for details). The incorporation of multiple memory cells allows for the mixing of information via the recurrent pathways $\mathbf{R}_{\mathbf{z}}$, $\mathbf{R}_{\mathbf{i}}$, $\mathbf{R}_{\mathbf{f}}$, $\mathbf{R}_{\mathbf{o}}$, which connect the hidden state $\mathbf{h}$ to the memory input $\mathbf{z}$ as well as to the gates $\mathbf{i}$, $\mathbf{f}$, $\mathbf{o}$. In this context, the use of exponential gating introduces a novel dimension to the process of memory mixing. Each sLSTM head is equipped to perform memory mixing internally, but no mixing occurs between different heads. The combination of head-based structuring and exponential gating in sLSTM gives rise to a fundamentally new mechanism for memory mixing.

\subsection{mLSTM}
\label{sec:mLSTM}

In order to improve the memory capabilities of LSTMs, we generalize the memory cell from a scalar $c \in \mathbb{R}$ to a matrix $\mathbf{C} \in \mathbb{R}^{d \times d}$. This reformulation allows for information retrieval via matrix multiplication. At each time step $t$, we aim to store a key vector $\mathbf{k}_t \in \mathbb{R}^d$ and a corresponding value vector $\mathbf{v}_t \in \mathbb{R}^d$, adopting the nomenclature commonly used in Transformers. Subsequently, at time $t + \tau$, a query vector $\mathbf{q}_{t+\tau} \in \mathbb{R}^d$ should allow access to the stored value $\mathbf{v}_t$. This framework aligns with the principles established in Bidirectional Associative Memories (BAMs) \citep{Kohonen:72,Anderson:72,Nakano:72,Anderson:77}.

To store each key-value pair, we employ the covariance update rule~\citep{Sejnowski:77,Dayan:91}

\begin{align}
\mathbf{C}_t \ &= \ \mathbf{C}_{t-1} \ + \ \mathbf{v}_{t} \ \mathbf{k}_t^\top \ .
\end{align}

We presuppose that a layer normalization is applied prior to mapping the inputs to keys and values, ensuring that these vectors have zero mean. The covariance update mechanism achieves optimality for maximizing the discriminability of the binary vectors being retrieved, which corresponds to optimizing the signal-to-noise ratio~\citep{Dayan:91}. Enhanced separability may be realized by restricting retrieval operations to pairwise associations, which, akin to attention mechanisms, incur quadratic computational costs~\citep{Krotov:16,Krotov:17,Ramsauer:21}. This covariance update procedure aligns with the Fast Weight Programmers paradigm~\citep{Schmidhuber:92ncfastweights,Schlag:21}, wherein it is common to employ a fixed decay factor scaling $\mathbf{C}_{t-1}$ and a fixed learning rate for the term $\mathbf{v}_{t} \mathbf{k}_t ^\top$~\citep{Ba:16}. Drawing on this connection, we embed the covariance update rule within the LSTM architecture: the forget gate implements the decay factor, the input gate serves as the learning rate, and the output gate modulates the amplitude of the retrieved representation.

In this matrix memory framework, the normalizer state is constructed as a weighted aggregation of key vectors, with each key vector scaled by both the corresponding input gate and the cumulative product of all subsequent forget gates. This configuration enables the normalizer state to track the magnitude of the gating signals over time. Given that the dot product between a query and the normalizer state may be nearly zero, the approach adopts the absolute value of this dot product and applies a minimum threshold (commonly set at 1.0), consistent with previous work~\citep{Sun:23arxiv}. The forward computation for the mLSTM is given as:

\begin{align}
\label{eq:mlstm_recurrent_begin}
\mathbf{C}_t \ &= \  f_t \ \mathbf{C}_{t-1} \ + \ 
  i_t \ \mathbf{v}_t \ \mathbf{k}_t^\top & &  &\text{cell state} \\
\mathbf{n}_t \ &= \  f_t \ \mathbf{n}_{t-1} \ + \ 
  i_t \ \mathbf{k}_t & &  &\text{normalizer state} \\
\mathbf{h}_t  \ &= \ \mathbf{o}_t \ \odot \ \tilde{\mathbf{h}}_t
\ , \qquad \qquad 
\tilde{\mathbf{h}}_t \ = \   \mathbf{C}_t \mathbf{q}_t \ / \ 
  \max \left\{ |\mathbf{n}_t^\top \mathbf{q}_t|, 1 \right\} 
   &&&\text{hidden state} \\
\mathbf{q}_t \ &= \ \mathbf{W}_q \ \mathbf{x}_t \ + \ \mathbf{b}_q  & &  &\text{query input} \\
\mathbf{k}_t \ &= \ \frac{1}{\sqrt{d}} \mathbf{W}_k \ \mathbf{x}_t \ + \ \mathbf{b}_k  & &  &\text{key input} \\
\mathbf{v}_t \ &= \ \mathbf{W}_v \ \mathbf{x}_t \ + \ \mathbf{b}_v  & &  &\text{value input} \\
i_t \ &= \ \!\exp \!  \! \left( \tilde{i}_t  \right) \ , 
 \qquad \qquad \ \ \ \ \,
  \tilde{i}_t \ = \ \mathbf{w}^\top_{i} \ \mathbf{x}_t \ + \  b_{i} \ \ \ 
  &&&\text{input gate} \\
f_t \ &= \ \sigma \!  \left(  \tilde{f}_t \right) \ \text{OR} \ \!\exp \!  \! \left(  \tilde{f}_t \right) , \ \ \: \!
\tilde{f}_t \ = \ \mathbf{w}^\top_{f} \ \mathbf{x}_t  \ + \
  b_{f}
  &&& \text{forget gate} \\
\label{eq:mlstm_recurrent_end}
\mathbf{o}_t \ &= \ \sigma \left( \tilde{\mathbf{o}}_t \right) \ , \qquad \qquad \quad \ \ \ \,
  \tilde{\mathbf{o}}_t  \ = \ \mathbf{W}_{\mathbf{o}} \ \mathbf{x}_t \ + \
  \mathbf{b}_{\mathbf{o}}
  &&&\text{output gate} 
\end{align}

The architecture of mLSTM supports the inclusion of several memory cells, similar to the standard LSTM model. In mLSTM, the distinction between having multiple heads and multiple cells is negligible due to the absence of memory mixing mechanisms. To ensure stability in the exponential gating functions within mLSTM, we adopt the stabilization strategies utilized for sLSTM, as described in Equation~\ref{eq:slstmstabil}. Owing to the lack of memory interaction in mLSTM, its recurrent operations can also be reformulated into a parallelizable format. Additional information can be found in Appendix~\ref{sec:appmLSTM}.

\subsection{xLSTM Architecture}
\label{sec:xLSTMarch}

\paragraph{xLSTM Blocks.}
\label{sec:xLSTMblock}
The xLSTM block is designed to capture and represent previous information in a non-linear, high-dimensional manner, thereby facilitating improved discrimination among various historical contexts. This enhanced separation of historical states is essential for the accurate prediction of subsequent elements in a sequence, such as the next token. The underlying principle is inspired by Cover’s Theorem~\citep{Cover:65}, which asserts that patterns subjected to non-linear transformations into higher dimensional spaces are more likely to become linearly separable relative to their original representations. We explore two forms of residual block architectures: (i) The first employs a residual block with a post up-projection step, akin to the approach in Transformers. Here, the model initially constructs a non-linear summary of the past in the original feature space, followed by a linear projection into a higher-dimensional space; a non-linear activation is then applied before projecting the result linearly back to the original space. This configuration is illustrated in the left section of Figure~\ref{fig:Backbones}
as well as in the third column of Figure~\ref{fig:xlstm_sketch}. A more granular illustration is available in Figure~\ref{fig:appsLSTM_detailed}
within the appendix. (ii) The second design features a residual block with a pre up-projection, following the design paradigm of State Space Models, where a linear transformation into a higher-dimensional space is performed at the outset,

The method first applies a nonlinear summary of the past within a high-dimensional space, followed by a linear projection back into the original dimensionality. In instances where an xLSTM block incorporates an sLSTM, the architecture utilizes the post up-projection block. Conversely, when an mLSTM is present within an xLSTM block, the pre up-projection block is adopted to accommodate the increased memory capacity afforded by the higher-dimensional space. For further clarification, readers are directed to the left panel of Figure~\ref{fig:Backbones}, the third column of Figure~\ref{fig:xlstm_sketch}, or Figure~\ref{fig:appsLSTM_detailed} in the appendix.

\paragraph{xLSTM Architecture. }
\label{sec:xLSTMarchitecure}
The xLSTM architecture is formed by stacking fundamental components in a residual manner~\citep{Srivastava:15,He:16}. Our implementation utilizes pre-LayerNorm~\citep{Ba:16b} residual backbones, a standard practice in state-of-the-art Large Language Models. This structure is illustrated in the final column of Figure~\ref{fig:xlstm_sketch}.

\subsection{Memory and Speed Considerations}

In contrast to Transformers, xLSTM networks exhibit linear computational costs and maintain fixed memory requirements regardless of sequence length. The compressive nature of xLSTM memory makes it particularly advantageous for industrial use cases and deployment on edge devices.

The mLSTM architecture achieves memory without relying on additional parameters; however, this comes at the cost of significant computational demands due to the $d \times d$ matrix-based memory and the corresponding $d \times d$ updates. We balance the size of the memory with the associated computational load. Despite this, since these operations can be parallelized effectively on GPUs, their impact on the overall computation time remains minimal.

Although mLSTM can be parallelized similarly to methods like FlashAttention~\citep{Dao:22,Dao:23} or GLA~\citep{Yang:23arxiv}, the sLSTM architecture cannot be parallelized owing to the presence of memory mixing through hidden-hidden connections. Nonetheless, we introduced an efficient CUDA-based implementation, incorporating GPU memory optimizations down to the register level, which typically results in only a less-than-twofold slowdown compared to mLSTM.

\section{Related Work}
\label{sec:related}

\paragraph{Linear Attention.} Numerous approaches have been proposed to address the Transformer's quadratic scaling with respect to context length by making attention computation linear in this parameter. The Synthesizer replaces traditional attention scores with synthetically learned weights, thereby avoiding direct token--token interactions~\citep{Tay:20arxiv}. Linformer achieves a linear self-attention mechanism through the use of low-rank projections, which enable efficient approximation~\citep{Wang:20arxiv}. Linear Transformer implements a linear version of attention computation~\citep{Katharopoulos:20}. Performer leverages a positive orthogonal random features technique to approximate the attention softmax linearly~\citep{Choromanski:21}. Additionally, alternatives such as Structured Global Convolution (SGConv)~\citep{Li:22} and Hyena Hierarchy~\citep{Poli:23} substitute conventional attention with rapid, large-kernel convolutional operations.

\paragraph{State Space Models.} In recent years, State Space Models (SSMs) have gained significant traction due to their linear scaling with sequence length and competitive results relative to Transformers. Early developments in this area include the Structured State Space sequence model (S4)~\citep{Gu:21}. This was succeeded by several variants such as the Diagonal State Space (DSS) model~\citep{Gupta:22}, Gated State Space (GSS) models~\citep{Mehta:22}, the S5 model~\citep{Smith:22}, Bidirectional Gated SSM (BiGS)~\citep{Wang:22}, H3 model~\citep{Fu:23}, and Mamba~\citep{Gu:24arxiv}.

\paragraph{Recurrent Neural Networks.} Recurrent Neural Networks (RNNs) have emerged as potential alternatives to Transformer architectures and attention mechanisms, primarily because of their linear scaling with respect to input sequence length. Various architectures, including Deep Linear Recurrent Units (LRUs), have demonstrated strong capabilities in language modeling tasks~\citep{Orvieto:23, De:24arxiv}. Hierarchically Gated Linear RNN (HGRN)~\citep{Qin:23} and its successor HGRN2~\citep{Qin:24arxiv} have also achieved notable results. Among the prominent RNN-based frameworks for large-scale language modeling is RWKV~\citep{Peng:23arxivshort,Peng:24arxivshort}, which rivals the performance of Transformer models.

\paragraph{Gating.}
Gating, a foundational concept introduced in LSTM architectures, has been rediscovered and reconceptualized across numerous recent methods. This mechanism appears in HGRN~\citep{Qin:23}, HGRN2~\citep{Qin:24arxiv}, Gated Linear Attention (GLA)~\citep{Yang:23arxiv}, Gated State Space (GSS) models~\citep{Mehta:22}, Bidirectional Gated SSM (BiGS)~\citep{Wang:22}, Moving Average Equipped Gated Attention (MEGA)~\citep{Ma:22}, RWKV~\citep{Peng:23arxivshort}, and Mamba~\citep{Gu:24arxiv}.

\paragraph{Covariance Update Rule.} In order to improve memory capacity, we incorporated a matrix memory into the mLSTM cell governed by a covariance update rule. Similar approaches employing this form of update include Fast Weight Programmers \citep{Schmidhuber:92ncfastweights,Schlag:21}, RWKV-5 and RWKV-6~\citep{Peng:24arxivshort}, Retention~\citep{Sun:23arxiv}, Linear Transformer~\citep{Katharopoulos:20}, and HGRN2~\citep{Qin:24arxiv}.

\paragraph{Most Related.} From a conceptual standpoint, the models most akin to xLSTM are Retention~\citep{Sun:23arxiv}, RWKV~\citep{Peng:23arxivshort,Peng:24arxivshort}, and HGRN2~\citep{Qin:24arxiv}. These architectures incorporate the ideas of matrix-based memory and/or employ gating mechanisms. Nevertheless, unlike the newly introduced sLSTM, these models do not facilitate memory mixing. The capacity for memory mixing is crucial for addressing state tracking challenges, making LSTMs inherently more expressive than State Space Models (SSMs) and Transformers~\citep{Merrill:24,Deletang:23}. Tasks such as code evaluation or maintaining awareness of entities throughout extensive narratives require robust state tracking capabilities.

\paragraph{Residually Stacking Architectures.} In alignment with most modern large-scale deep learning systems, xLSTM architectures employ a strategy of assembling components through residual stacking~\citep{Srivastava:15,He:16}.

This architectural design has been fundamental in enabling the development of deep convolutional neural networks~\citep{He:16} and is integral to the architecture of Transformers~\citep{Vaswani:17}. In fact, Transformers underpin the recent surge in Large Language Models (LLMs) such as GPT-3~\citep{Brown:20short}, ChatGPT~\citep{Schulman:22short}, GPT-4~\citep{Achiam:23gpt4short}, Megatron-LM~\citep{Shoeybi:19}, Gopher~\citep{Rae:21short}, ERNIE 3.0 Titan~\citep{Wang:21arxivshort}, GLaM~\citep{Du:21glamshort}, Chinese M6~\citep{Lin:21}, the multilingual AlexaTM 20B~\citep{Soltan:22}, OPT~\citep{Zhang:22opt}, Chinchilla~\citep{Hoffmann:22short}, BLOOM~\citep{Scao:22arxivshort}, GLM-130B~\citep{Zeng:22arxivshort}, LaMDA~\citep{Thoppilan:22short}, PaLM~\citep{Chowdhery:22short}, Llama~\citep{Touvron:23arxiv}, and Gemini~\citep{GeminiTeam:23arxiv,Reid:24arxiv}.

\section{Experiments}
\label{sec:Exp}

In this section, we present an empirical assessment of xLSTM, juxtaposing its performance with that of established techniques, particularly in the domain of language modeling. Section~\ref{sec:ExpSyntethic} explores xLSTM’s distinct abilities by applying it to carefully constructed synthetic tasks. Subsequently, in Section~\ref{sec:ExpComparison}, we compare the validation perplexity scores achieved by a range of recent language modeling approaches trained using 15B tokens from SlimPajama~\citep{Soboleva:23}. We also perform ablation analyses of xLSTM on the same data. Additionally, the scaling characteristics of each approach are examined in accordance with methodologies proposed by \citet{Kaplan:20} and \citet{Brown:20short}.

Section~\ref{sec:ExpLanguage} extends our evaluation with a comprehensive language modeling study. Here, xLSTM and the leading methods identified in Section~\ref{sec:ExpComparison} are trained on a substantially larger set of 300B tokens from SlimPajama~\citep{Soboleva:23}. Our investigation proceeds along several axes: (1) we test the ability of each approach to extrapolate effectively to substantially longer context windows; (2) we compare their validation perplexity and results on a suite of downstream tasks~\citep{Sutawika:24}; (3) performance is further measured across 571 text domains included in the PALOMA language benchmark dataset~\citep{Magnusson:23arxivshort}; and (4) we revisit the scaling analysis of all methods, now with a twentyfold increase in the volume of training data.

Throughout our experiments, we denote xLSTM[$a$:$b$] as representing the ratio $a/b$ of mLSTM-based to sLSTM-based xLSTM blocks. For instance, xLSTM[7:1] indicates that, among eight blocks, seven utilize mLSTM while one employs sLSTM. When the total number of blocks is 48—a typical setting—this configuration corresponds to 42 mLSTM-based blocks and 6 sLSTM-based blocks. Additionally, each experiment incorporates pre and post up-projection blocks for mLSTM and sLSTM, respectively.

\subsection{Synthetic Tasks and Long Range Arena}
\label{sec:ExpSyntethic}
To begin, we evaluate the impact of xLSTM’s exponential gating combined with memory mixing on formal language benchmarks~\citep{Deletang:23}. Next, we investigate the capabilities of xLSTM's matrix memory in addressing the Multi-Query Associative Recall challenge~\citep{Arora:23arxiv}. Lastly, we measure how well xLSTM handles extended sequences using the Long Range Arena suite~\citep{Tay:21}.

\paragraph{Assessment of xLSTM's Exponential Gating Mechanism with Memory Mixing.} We evaluate the newly introduced exponential gating combined with memory mixing in xLSTM, a mechanism designed to facilitate state tracking tasks~\citep{Merrill:24, Merrill:23}. The formal language benchmarks developed by \citet{Deletang:23} are adapted and further expanded to allow training on varying sequence lengths, thereby supporting extrapolation to longer sequences. Comprehensive descriptions of the individual tasks along with supplementary results can be found in Appendix~\ref{sec:appExpSynthetic-formalLang}. Our comparative analysis includes xLSTM and other established architectures, such as Transformers, State Space Models, and traditional Recurrent Neural Networks. For evaluation, we measure accuracy on task-specific relevant tokens, normalizing scores such that 0 represents chance level and 1 denotes flawless performance. We focus on 2-block model configurations from each approach, specifically: xLSTM[0:1] (sLSTM only), xLSTM[1:0] (mLSTM only), xLSTM[1:1], Llama, Mamba, RWKV, Retention, Hyena, LSTM, and a variant of LSTM incorporated within Transformer blocks (LSTM (Block)). The performance outcomes for these methods are depicted in Figure~\ref{fig:formal-main}. Notably, architectures such as Transformers and State Space Models, in the absence of memory mixing (i.e., lacking state-tracking capabilities), fail to successfully handle certain regular grammars, exemplified by the parity task. This observation corroborates earlier work indicating that Transformers and State Space Models are inherently less capable than RNNs in such scenarios~\citep{Merrill:24, Merrill:23, Deletang:23}.

\paragraph{Assessment of xLSTM's Memory Capacity via Associative Recall Tasks.} In this study, we evaluate the matrix memory of xLSTM by measuring its capacity on the Multi-Query Associative Recall benchmark~\citep{Arora:23arxiv}. In each input sequence, randomly selected key-value pairs from an extensive vocabulary must be stored and retrieved at a later stage. To increase the challenge beyond that posed by the original benchmark, the number of key-value pairs is raised up to 256 and the context length is extended to 2048, providing a comprehensive evaluation of the memory capacity across different model architectures. We conduct comparisons across 2-block variants of Llama, Mamba, RWKV-5, RWKV-6, xLSTM[1:1], and xLSTM[1:0], gauging performance according to recall accuracy of key-value associations. Because Transformer models such as Llama exhibit memory capabilities that scale exponentially with their representation dimension~\citep{Ramsauer:21}, they serve as the performance upper bound for this type of memory test. Figure~\ref{fig:mqar-main} summarizes these outcomes. xLSTM[1:1] demonstrates superior memory performance among non-Transformer models, even at reduced model sizes. Notably, inclusion of the sLSTM block does not compromise memory but instead strengthens it, especially at the highest difficulty level with 256 pairs. Supplementary analyses and detailed results appear in Appendix~\ref{sec:appExpSynthetic-mqar}, where extrapolation studies suggest that the advanced memory of xLSTM facilitates generalization to context lengths surpassing those encountered during training.

\paragraph{Evaluation of xLSTM's Proficiency with Extended Contexts Using Long Range Arena.} In order to evaluate how effectively xLSTM manages lengthy input sequences and expansive contexts, we benchmark various approaches on the Long Range Arena~\citep{Tay:21}. xLSTM achieves robust and reliable results across all tested tasks, indicating that the architecture possesses a high degree of effectiveness when it comes to managing diverse challenges associated with extended contexts.

Further information is provided in Appendix~\ref{sec:appExpSynthetic-lra}.

\subsection{Method Comparison and Ablation Study}
\label{sec:ExpComparison}

In order to explore the central inquiry of this work—namely, the capabilities of our novel LSTM variants when scaled for language modelling—we conduct experiments by training xLSTMs, Transformers, State Space Models, as well as alternative approaches, utilizing 15B tokens from SlimPajama under a unified auto-regressive framework. We evaluate all trained models on the validation set and carry out ablation experiments specifically targeting the xLSTM architectures.

\paragraph{Comparing xLSTM to Other Methods.} To facilitate a direct comparison, we trained various models using a corpus of 15B tokens sourced from SlimPajama~\citep{Soboleva:23}. Model performance is assessed by measuring perplexity on a designated validation set. The evaluation includes a diverse selection of models: our proposed xLSTM, GPT-3 (Transformer)~\citep{Brown:20short}, Llama (Transformer)~\citep{Touvron:23arxiv}, H3 (SSM)~\citep{Fu:23}, Mamba (SSM)~\citep{Gu:24arxiv}, RWKV-4, RWKV-5, and RWKV-6 (all RNNs)~\citep{Peng:23arxivshort, Peng:24arxivshort}, GLA (a linear Transformer)~\citep{Yang:23arxiv}, HGRN and HGRN2 (RNNs)~\citep{Qin:23, Qin:24arxiv}, RetNet (linear Transformer)~\citep{Sun:23arxiv}, and Hyena (linear Transformer)~\citep{Poli:23}, in addition to xLSTM[1:0] and xLSTM[7:1] as defined in Section~\ref{sec:xlstm_blocks_notation}. Training across models was conducted using mixed precision; however, for RWKV-5, RWKV-6, GLA, and HGRN2, mixed-precision training was implemented independently of PyTorch’s automated tools (see also Appendix Section~\ref{sec:appGenTrainProc}). The examined models are grouped as follows: (a) Transformer architectures, (b) State Space Models (SSMs), and (c) Recurrent Neural Networks (RNNs) as well as linear Transformers, which employ a linear mechanism in place of conventional attention in Transformers. All models are calibrated to be equivalent in scale to a GPT-3 model containing 350M parameters, featuring an embedding dimension of 1024 with 24 residual blocks. Notably, only GPT-3 utilizes shared weights between token and output embeddings, resulting in a reduced parameter count. As reported in Table~\ref{tab:spaj15b_model_results}, xLSTM achieves the lowest validation perplexity among all compared approaches. Additional information is provided in Appendix~\ref{sec:appExpComparison}. Figure~\ref{fig:temp_sclaw15B} depicts the scaling trends observed in this experiment, suggesting that xLSTM’s superior performance is likely to persist as model size increases.

\paragraph{Ablation Studies.}
As illustrated in Table~\ref{tab:spaj15b_model_results} and Figure~\ref{fig:temp_sclaw15B}, xLSTM delivers strong performance in language modeling tasks when trained on 15B tokens from SlimPajama. To isolate the contributions of each architectural change, we systematically transform a standard LSTM into an xLSTM in a sequence of modifications. The process begins by embedding LSTM layers within pre-LayerNorm residual structures. The next modification introduces a post up-projection block. The final stage incorporates exponential gating and matrix memory components.

Table~\ref{tab:ablstudies} (top) presents the experimental outcomes.

The ablation analysis identifies both the exponential gating and the matrix memory as critical factors driving the observed gains in performance. Moreover, given the centrality of gating in RNNs and State Space Models, we evaluate the effects of various gating strategies. From the results in Table~\ref{tab:ablstudies} (bottom), we find that learnable gates, particularly those modulated by input, consistently enhance performance. Further investigation into the specific backbone components is provided in Appendix~\ref{sec:appAblDetails}.

\subsection{xLSTM as Large Language Model}
\label{sec:ExpLanguage}

This work concludes with extensive language modeling experiments to evaluate the viability of xLSTM as a large language model (LLM). To this end, we expand the dataset and train using 300B tokens from SlimPajama. Notably, this token count aligns with the training regimes of models such as Mamba~\citep{Gu:24arxiv} and Griffin~\citep{De:24arxiv}.

We benchmark xLSTM against representative models from each approach outlined in Section~\ref{sec:ExpComparison}, specifically: RWKV-4, Llama, and Mamba. RWKV-4 is chosen to represent the RNN family because stable and effective training precision parameters for RWKV-5, RWKV-6, and HGRN2 were not identified until after the initiation of the 300B token training runs (refer to Appendix~\ref{sec:appGenTrainProc} for further details). 

Our experimental setup involves training various model sizes (125M, 350M, 760M, 1.3B). Each model undergoes evaluation for its ability to extrapolate to longer sequences, alongside performance assessments on the validation dataset. Further, we evaluate the models' downstream task competencies, their effectiveness in language modeling across 571 distinct text domains as delineated in the PALOMA benchmark, and analyze how their performance scales with increased capacity.

\paragraph{Sequence Length Extrapolation.}
\label{sec:spaj300B_ctx_extrapolation}
We begin by evaluating the ability of 1.3B parameter models—including xLSTM, RWKV-4, Llama, and Mamba—to extrapolate to longer sequence lengths. Each model is trained with a fixed context length of 2048 tokens and subsequently evaluated on extended context lengths up to 16384 tokens. The performance metrics are presented in Figure~\ref{fig:spaj300B_seq_extrapolation}. Notably, xLSTM consistently demonstrates lower perplexity on longer contexts compared to the other architectures.

\paragraph{Validation Perplexity and Downstream Tasks.}
\label{sec:spaj_downstream_eval}
Next, across every model scale, we assess how xLSTM, RWKV-4, Llama, and Mamba perform on the SlimPajama validation dataset for next token prediction. Additionally, we measure their effectiveness on a suite of downstream tasks that probe common sense reasoning abilities.

The third column in Table~\ref{tab:lmeval} presents the perplexity scores on the validation set for various approaches. Both xLSTM[1:0] and xLSTM[7:1] achieve the lowest validation set perplexity across all considered model scales. Performance results for downstream tasks are shown in the remaining columns of Table~\ref{tab:lmeval}. Across nearly all downstream tasks and model sizes, xLSTM consistently outperforms other methods, except for the ARC task, where Mamba occasionally yields the top results. Additional information is available in Appendix~\ref{sec:appExpLanguage}.

\paragraph{Performance on PALOMA Language Tasks.} Thirdly, we evaluate the next token prediction capabilities of xLSTM, RWKV-4, Llama, and Mamba models, across all parameter sizes, using PALOMA language tasks~\citep{Magnusson:23arxivshort}. Specifically, we assess model performance based on perplexity scores for next token prediction across 571 distinct text domains, including sources such as nytimes.com and Reddit's r/depression community. Table~\ref{tab:ppleval} presents perplexity values, organized by language modeling benchmarks (the first seven columns) and by specific domain evaluations (the final five columns). On these tasks, xLSTM[1:0] consistently outperforms xLSTM[7:1].

xLSTM[1:0] achieves lower perplexity compared to Mamba on 568 of 571 (99.5\%) text domains, outperforms Llama in terms of perplexity on 486 of 571 (85.1\%) domains, and surpasses RWKV-4 with reduced perplexity on 570 of 571 (99.8\%) domains. Additional information is available in Appendix~\ref{sec:appExpLanguage}.

\paragraph{Scaling Laws.} Fourth, we investigate the presence of power-law scaling, which enables prediction of model performance at larger scales~\citep{Kaplan:20,Brown:20short}. The scaling trends are visualized in Figure~\ref{fig:temp_sclaw300B}. While all architectures exhibit analogous scaling patterns, their baseline performance levels differ. Among them, RWKV-4 demonstrates the lowest performance, succeeded by Llama and Mamba. xLSTM surpasses Mamba, exhibiting an advantage over Mamba similar to that which Mamba has over Llama. These results suggest that xLSTM is likely to retain its performance advantage over both Transformers and State-Space models as the model size increases.

\textbf{Generation Times and Maximal Throughput.} In Figure~\ref{fig:inference_time_generation}, we report the time required for text generation, while Figure~\ref{fig:inference_time_decoding_throughput} (left) presents the peak throughput for various xLSTM configurations at the 1.3B parameter level. Our evaluation also includes similarly sized Mamba, Llama, and RWKV implementations sourced from HuggingFace, and for Llama, we incorporate a static key-value cache. At the point when these experiments were carried out, it was not possible to utilize full cache compilation for the Transformer Model or employ \texttt{torch.compile} for Mamba model compilation. During the text generation benchmarks, we use a batch size of 1 for all models, with a pre-fill of 16 tokens, which is particularly advantageous for the Transformer architecture.

As illustrated in Figure~\ref{fig:inference_time_generation}, xLSTM, together with the other recurrent architectures Mamba and RWKV-4, demonstrates linear computational scaling, whereas Llama exhibits quadratic scaling behavior. For the analysis of decoding throughput, we assess a variety of batch sizes and prefill settings for the Llama model.

Figure~\ref{fig:inference_time_decoding_throughput} (right) reveals that xLSTM supports significantly larger batch sizes in comparison to Llama, attributed to its constant memory footprint, enabling xLSTM to attain superior throughput rates.

\section{Limitations}

(i) Unlike mLSTM, the sLSTM’s approach to memory mixing hinders the use of parallelized computations, precluding efficient parallel implementations. Despite this, we managed to create a custom CUDA kernel for sLSTM, achieving execution speeds less than twice as slow as our parallel mLSTM version. (ii) Presently, the CUDA kernels for mLSTM lack optimization, rendering the current setup approximately four times slower than either FlashAttention or the scan utilized in Mamba. Enhanced CUDA kernel performance is likely attainable by adopting strategies similar to those employed in FlashAttention. (iii) The matrix-based memory architecture of mLSTM incurs significant computational complexity, as it requires manipulation of $d \times d$ matrices. However, since both the update and readout operations are parameter-free and can leverage standard matrix computation parallelism, the actual wall-clock performance penalty introduced by this memory design is relatively modest. (iv) Appropriate initialization of the forget gates is essential for proper operation. (v) The memory matrix in mLSTM operates independently of sequence length, which poses a risk of memory saturation when handling longer sequences. Nevertheless, this has not manifested as a bottleneck for contexts of up to 16k tokens; see Section~\ref{sec:spaj300B_ctx_extrapolation}. (vi) Given the considerable computational cost associated with training large-scale language models, we have not exhaustively tuned either the architectural design or hyperparameters for xLSTM, particularly in the case of more expansive configurations. We expect that rigorous and thorough optimization will be required for xLSTM architectures to realize their optimal performance.

\section{Conclusion}
We have provided a partial answer to our guiding question: To what extent can language modeling be advanced by scaling LSTM architectures to billions of parameters? Our findings show that such scaling allows LSTM-based models to achieve performance levels at least on par with prevailing approaches, including Transformers and State Space Models. Through the introduction of exponential gating, memory mixing, and an innovative memory architecture, we have extended the capabilities of the LSTM, resulting in the xLSTM. Empirically, xLSTM demonstrates strong results in language modeling, frequently outperforming leading alternatives such as Transformers and State Space Models. Observed scaling behavior suggests that even larger xLSTM models could emerge as formidable challengers to Transformer-based Large Language Models. Furthermore, xLSTM holds promise for transformative applications in areas beyond language modeling, such as Reinforcement Learning, Time Series Prediction, and the modeling of physical processes.

\end{document}